\section*{Chapter \ref{ch:g9:algorithms-on-numbers} --- Algorithms on Numbers}

\begin{solution}{exo:g9:algorithms-on-numbers:1}
\code{[1, 2, 4, 7, 14, 28]}. The divisors smaller than 28 add up to
$1 + 2 + 4 + 7 + 14 = 28$: the number itself. 28 is a perfect number
(Exercise \ref{exo:g9:algorithms-on-numbers:4}).
\end{solution}

\begin{solution}{exo:g9:algorithms-on-numbers:2}
Left: 2, 3, 5, 7, 11, 13, 17, 19, 23, 25, 29. All are prime except 25, a
multiple of 5: the sieve crosses it out at the next turn, when 5 is taken.
\end{solution}

\begin{solution}{exo:g9:algorithms-on-numbers:3}
$(48, 18) \to (18, 12) \to (12, 6) \to (6, 0)$: three steps, and the result is 6.
\end{solution}

\begin{solution}{exo:g9:algorithms-on-numbers:4}
\omcode{code/grade-9/04-algorithms-on-numbers/exercises/perfect.py}{1}{12}{Perfect
numbers below 1,000.}\label{lst:g9:algorithms-on-numbers:sol-perfect}

It prints 6, 28 and 496. A divisor smaller than $n$ is at most $n / 2$, so the
\omterm{def:g4:algorithms-repeat-and-decide:loop}{loop} stops there.
\end{solution}

\begin{solution}{exo:g9:algorithms-on-numbers:5}
For $n = 1$, the \omterm{def:g4:algorithms-repeat-and-decide:condition}{condition} \code{2 * 2 <= 1} is false at once: the \omterm{def:g4:algorithms-repeat-and-decide:loop}{loop} never
runs, and the \omterm{def:g7:functions:function}{function} reaches \code{return True}. The fix rules out the numbers
below 2 first:

\omcode{code/grade-9/04-algorithms-on-numbers/exercises/is_prime_one_fixed.py}{1}{9}{Numbers
below 2 are not prime.}\label{lst:g9:algorithms-on-numbers:sol-one}
\end{solution}

\begin{solution}{exo:g9:algorithms-on-numbers:6}
The squares of primes: 4, 9, 25, 49, and so on. For 4, the \omterm{def:g4:algorithms-repeat-and-decide:loop}{loop} needs
\code{2 * 2 < 4}, which is false, so 2 is never tried; the smallest wrong answer
is 4. A square of a prime has no divisor below its square root, only the root
itself, so the \omterm{def:g4:algorithms-repeat-and-decide:loop}{loop} must include $d \times d = n$: write \code{d * d <= n}.
\end{solution}

\begin{solution}{exo:g9:algorithms-on-numbers:7}
\omcode{code/grade-9/04-algorithms-on-numbers/exercises/gcd3.py}{1}{12}{The gcd of
three numbers.}\label{lst:g9:algorithms-on-numbers:sol-gcd3}

It prints 42. A number divides $a$, $b$ and $c$ exactly when it divides
$\gcd(a, b)$ and $c$.
\end{solution}

\begin{solution}{exo:g9:algorithms-on-numbers:8}
\omcode{code/grade-9/04-algorithms-on-numbers/exercises/twins.py}{1}{16}{Twin
primes below 1,000.}\label{lst:g9:algorithms-on-numbers:sol-twins}

It prints \code{35 pairs; the last: [881, 883]}.
\end{solution}

\begin{solution}{exo:g9:algorithms-on-numbers:9}
A multiple $k \times p$ with $k < p$ has a factor $k$ smaller than $p$, so a
prime smaller than $p$ divides it. That prime was taken earlier, and its
multiples, $k \times p$ among them, are already crossed out.
\end{solution}

\begin{solution}{exo:g9:algorithms-on-numbers:10}
$100 \times 100 = 10{,}000 \le 10{,}007$ but $101 \times 101 = 10{,}201$: the \omterm{def:g4:algorithms-repeat-and-decide:loop}{loop}
tries $d = 2$ to 100, that is 99 divisions, and finds no divisor. A test of the
chapter checks it.
\end{solution}

\begin{solution}{exo:g9:algorithms-on-numbers:11}
Every two steps, the larger number falls below half of what it was. A number
below 1,000 is below $2^{10} = 1{,}024$, so after 20 steps the larger number
would be below 1, which cannot be: the \omterm{def:g4:algorithms-repeat-and-decide:algorithm}{algorithm} has stopped by then. At most
about 20 steps; the real record is 14, for $(987, 610)$.
\end{solution}

\begin{solution}{exo:g9:algorithms-on-numbers:12}
Saved next to \code{is_prime.py}, the \omterm{def:g7:functions:function}{function} tries each candidate in turn
with the fast test, adding up the divisions, while the candidate is not prime:

\omcode{code/grade-9/04-algorithms-on-numbers/exercises/next_prime.py}{1}{17}{The
next prime and its cost.}\label{lst:g9:algorithms-on-numbers:sol-next}

For 1,000 it returns \code{(1009, 76)}: 1,001 to 1,008 are thrown out after a
few divisions each, and 1,009 needs 30. A test checks these \omterm{def:g6:variables-and-input:value}{values}.
\end{solution}

\begin{solution}{pb:g9:algorithms-on-numbers:1}
\textbf{1.} $(13, 8) \to (8, 5) \to (5, 3) \to (3, 2) \to (2, 1) \to (1, 0)$:
5 steps.

\textbf{2.} One more, 6 steps: the six cuts of the figure.

\textbf{3.} One step: $100 \bmod 1 = 0$. Size alone does not make a pair slow.

\textbf{4.} Every pair with $1 \le b < a \le 999$: $999 \times 998 / 2 =
498{,}501$ pairs.

\textbf{5.} It prints a pair only when it beats the best count so far.

\omoutput{code/grade-9/04-algorithms-on-numbers/outputs/slowest_pairs.txt}{The record
pairs below 1,000.}\label{out:g9:algorithms-on-numbers:slowest}

\textbf{6.} $(987, 610)$, with 14 steps.

\textbf{7.} Each number is the sum of the two before it: $13 = 8 + 5$,
$21 = 13 + 8$, \ldots: the sequence 1, 1, 2, 3, 5, 8, 13, \ldots

\textbf{8.} $(987 + 610, 987) = (1597, 987)$, with 15 steps.

\textbf{9.} 1 at every step, except the last one, $(2, 1)$, where it is 2.

\textbf{10.} If $a = b + c$ for the number $c$ just before $b$ in the sequence,
then $c < b$, so $a \bmod b = c$: the step turns $(a, b)$ into $(b, c)$, the
record pair just before.

\textbf{11.} A quotient of 1 removes as little as possible from $a$. Larger
quotients make the numbers fall faster, so fewer steps are needed.

\textbf{12.} \textbf{Result.} Below 1,000, the pair $(987, 610)$ takes the most
division steps, 14. The record pairs are neighbouring numbers of the sequence
1, 1, 2, 3, 5, 8, 13, \ldots, because each step of Euclid takes such a pair
back exactly one place.
\end{solution}
